%! Setting Up the SVD — Unit 7, Lesson 7.0 — Slides (Beamer)
\documentclass[aspectratio=169, 11pt]{beamer}
\usepackage{linalg-beamer}   % provides \burgundyheader{} and \sectionlabel[color]{}
\newcommand{\uu}{\mathbf{u}}
\newcommand{\vv}{\mathbf{v}}
\newcommand{\ee}{\mathbf{e}}
\newcommand{\T}{^{\top}}

\begin{document}

% TITLE SLIDE (CourseName is not defined in beamer, so the name is literal)
{
\setbeamercolor{background canvas}{bg=burgundy}
\begin{frame}[plain]
  \vspace{2.8cm}\hspace{0.55cm}%
  \begin{minipage}{0.9\textwidth}
    {\color{goldacc}\bfseries\small LESSON 7.0}\\[0.5cm]
    {\color{white}\bfseries\LARGE Setting Up the SVD}\\[1.0cm]
    {\color{blushmid}\small Linear Algebra~$\cdot$~Unit 7: The Singular Value Decomposition}
  \end{minipage}
\end{frame}
}

% HOOK
\begin{frame}[t, plain]
  \burgundyheader{What does a matrix do to the unit circle?}
  \sectionlabel{Hook}
  \begin{columns}[T]
    \begin{column}{0.46\textwidth}
      In Unit~5 a matrix \textbf{moved} the plane; in Unit~6 we found the directions it only
      \textbf{stretched}. Now multiply every vector of length $1$ (the \emph{unit circle}) by a matrix:\\[6pt]
      the circle becomes an \textbf{ellipse}. Its two half-widths are the \textbf{singular values}
      $\sigma_1\ge\sigma_2\ge0$.
    \end{column}
    \begin{column}{0.52\textwidth}
      \centering
      \begin{tikzpicture}[scale=0.42]
        \begin{scope}[shift={(-4.4,0)}]
          \draw[->, charcoal, thin] (-1.8,0)--(1.8,0);
          \draw[->, charcoal, thin] (0,-1.8)--(0,1.8);
          \draw[burgundy!70, thick] (0,0) circle (1);
          \draw[->, ultra thick, burgundy] (0,0)--(0.707,0.707) node[above right=-2pt]{\scriptsize $\vv_1$};
          \draw[->, ultra thick, royalblue] (0,0)--(0.707,-0.707) node[below right=-2pt]{\scriptsize $\vv_2$};
        \end{scope}
        \draw[->, very thick, goldacc] (-2.0,0)--(-0.2,0) node[midway, above]{\scriptsize $A$};
        \begin{scope}[shift={(3.1,0)}]
          \draw[->, charcoal, thin] (-3.4,0)--(3.4,0);
          \draw[->, charcoal, thin] (0,-3.4)--(0,3.4);
          \draw[burgundy!70, thick, rotate=45] (0,0) ellipse (3 and 1);
          \draw[->, ultra thick, burgundy] (0,0)--(2.121,2.121) node[above right=-2pt]{\scriptsize $A\vv_1=3\uu_1$};
          \draw[->, ultra thick, royalblue] (0,0)--(0.707,-0.707) node[below right=-2pt]{\scriptsize $A\vv_2=\uu_2$};
        \end{scope}
      \end{tikzpicture}
    \end{column}
  \end{columns}
  \vspace{2pt}
  \begin{block}{The question of Unit 7}
    Which perpendicular directions does a matrix stretch, and \textbf{by how much}?
  \end{block}
\end{frame}

% DEFINITION
\begin{frame}[t, plain]
  \burgundyheader{Singular values and singular vectors}
  \sectionlabel[goldacc]{Definition}
  \begin{itemize}
    \setlength{\itemsep}{6pt}
    \item Every matrix sends a perpendicular \textbf{input} pair to a perpendicular \textbf{output} pair:
          \[ A\vv_1=\sigma_1\uu_1,\qquad A\vv_2=\sigma_2\uu_2,\qquad \sigma_1\ge\sigma_2\ge0. \]
    \item The $\vv$'s are the \textbf{input} singular vectors, the $\uu$'s the \textbf{output} singular
          vectors, the $\sigma$'s the \textbf{singular values} (the ellipse's half-widths).
    \item \textbf{Symmetric $A=\begin{bsmallmatrix} 2 & 1\\ 1 & 2 \end{bsmallmatrix}$:}
          $A\begin{bsmallmatrix} 1\\1 \end{bsmallmatrix}=3\begin{bsmallmatrix} 1\\1 \end{bsmallmatrix}$
          ($\sigma_1=3$), \quad
          $A\begin{bsmallmatrix} 1\\-1 \end{bsmallmatrix}=1\begin{bsmallmatrix} 1\\-1 \end{bsmallmatrix}$
          ($\sigma_2=1$).
    \item Singular values are \textbf{lengths}: always real and $\ge0$.
  \end{itemize}
\end{frame}

% READ IT OFF
\begin{frame}[t, plain]
  \burgundyheader{Reading it off --- input frame need not equal output frame}
  \sectionlabel{Skill}
  \begin{block}{The read-off}
    Feed a matrix a perpendicular pair; the outputs are perpendicular too. The output \emph{lengths} are
    the singular values (largest first); the output \emph{directions} are the $\uu$'s.
  \end{block}
  \begin{columns}[T]
    \begin{column}{0.5\textwidth}
      \textbf{Swap-stretch} $S=\begin{bsmallmatrix} 0 & 2\\ 1 & 0 \end{bsmallmatrix}$:
      \[ S\ee_1=\begin{bsmallmatrix} 0\\1 \end{bsmallmatrix},\quad S\ee_2=\begin{bsmallmatrix} 2\\0 \end{bsmallmatrix}. \]
      Perpendicular out; $\sigma_1=2,\ \sigma_2=1$.
    \end{column}
    \begin{column}{0.5\textwidth}
      Here $\vv_1=\ee_2$ is sent to $\uu_1=\begin{bsmallmatrix} 1\\0 \end{bsmallmatrix}$ --- a
      \textbf{different} direction. Input and output frames can differ.
    \end{column}
  \end{columns}
\end{frame}

% WHY + ROAD AHEAD
\begin{frame}[t, plain]
  \burgundyheader{Why it matters --- and the road ahead}
  \sectionlabel[goldacc]{Payoff}
  \begin{itemize}
    \setlength{\itemsep}{6pt}
    \item \textbf{The SVD:} collect the pieces $\Rightarrow$ $A=U\Sigma V\T$ --- rotate ($V\T$), stretch
          ($\Sigma$), rotate ($U$).
    \item \textbf{Works for every matrix} (even non-square), and $\sigma\ge0$ always --- unlike
          eigenvectors, which need a square matrix and may not exist.
    \item \textbf{Finding them next:} $\sigma=\sqrt{\lambda}$ for the eigenvalues $\lambda$ of
          \[ A\T A \quad\text{(a symmetric matrix --- Unit 6).} \]
  \end{itemize}
  \vspace{2pt}
  \begin{block}{The road ahead}
    \textbf{7.1} finding singular values/vectors \;$\bullet$\; \textbf{7.2} compressing images \;$\bullet$\;
    \textbf{7.3} principal component analysis \;$\bullet$\; \textbf{7.4} the victory of orthogonality.
  \end{block}
\end{frame}

\end{document}
